\begin{notechapter}
\chapter{The Geometry of the Gradient}\label{ch:gradient-tangent-space}\label{ch:gradient-misconception}
\chaptersubtitle{Where the gradient lives, what it depends on, and why common pictures mislead}

\begin{chapterintro}
Gradient descent, $\theta\leftarrow\theta - \eta\nabla L(\theta)$, is the workhorse of machine learning, and
its geometry is often misrepresented: the gradient is drawn as an arrow sliding down the loss
surface. This chapter sets the picture straight and shows that the gradient depends on a
choice of inner product --- which is exactly what preconditioning, Newton's method and the
natural gradient change. You will learn
\begin{itemize}[nosep]
  \item the difference between the differential (a linear form) and the gradient (a vector),
        and how an inner product turns one into the other;
  \item why the gradient is the direction of steepest ascent and is orthogonal to level sets;
  \item why the gradient lives in parameter space and not on the graph of the loss;
  \item which structure the update $\theta - \eta\nabla L(\theta)$ really uses (an affine space with an inner product),
        and how to detect formulas that secretly depend on the origin, such as weight decay;
  \item how changing the metric changes the gradient: preconditioning, Newton, natural gradient;
  \item gradients on curved parameter spaces (manifolds), with the example of the sphere and the
        link to PCA.
\end{itemize}
\prerequisites{multivariable calculus; \cref{ch:adjoint-operator} (inner products, Riesz,
projections); \cref{ch:pca-overview} for the last example.}
\end{chapterintro}

\section{Differential and gradient}\label{gradient-tangent-space:sec:differential}

Let $f:\R^n\to\R$ be differentiable at $\theta$. Its \emph{differential} is the linear map
$df_\theta:\R^n\to\R$ that best approximates $f$ near $\theta$:
\[
  f(\theta + v) = f(\theta) + df_\theta(v) + o(\norm v),\qquad
  df_\theta(v) = \sum_{k=1}^n\frac{\partial f}{\partial\theta_k}(\theta)\,v_k
  = \lim_{t\to0}\frac{f(\theta + tv) - f(\theta)}{t} .
\]
$df_\theta(v)$ is the \emph{directional derivative} of $f$ in the direction $v$. The differential is a
\emph{linear form} --- an element of the dual space, a ``row vector'' --- and it is defined without any
inner product. To get a \emph{vector} that can be added to $\theta$, we need an inner product and the
Riesz theorem (\cref{adjoint-operator:thm:riesz}).

\begin{definition}[Gradient with respect to an inner product]
Given an inner product $\inner{\cdot}{\cdot}$ on $\R^n$, the gradient $\nabla f(\theta)$ is the unique vector
such that
\[
  \inner{\nabla f(\theta)}{v} = df_\theta(v)\qquad\text{for all }v\in\R^n .
\]
\end{definition}

\begin{proposition}[Gradient for the metric $M$]\label{gradient-tangent-space:prop:metric}
For the dot product, $\nabla f(\theta) = \bigl(\frac{\partial f}{\partial\theta_1},\dots,\frac{\partial f}{\partial\theta_n}\bigr)^\top$, the vector of
partial derivatives. For $\inner uv_M = u^\top Mv$, the gradient is $\nabla_Mf(\theta) = M^{-1}\nabla f(\theta)$.
\end{proposition}

\begin{proof}
$df_\theta(v) = \sum_k\partial_kf\,v_k = (\nabla f)^\top v$, which is the defining identity for the dot
product. For $M$: we need $u^\top Mv = (\nabla f)^\top v$ for all $v$, i.e.\ $Mu = \nabla f$, so $u = M^{-1}\nabla f$.
\end{proof}

\begin{keyidea}
The partial derivatives define the differential, which is intrinsic. The \emph{gradient vector}
depends on how we measure lengths and angles. ``The'' gradient of calculus is the gradient for
the Euclidean dot product in the chosen coordinates.
\end{keyidea}

\begin{theorem}[Steepest ascent]\label{gradient-tangent-space:thm:steepest}
If $\nabla f(\theta)\neq0$, then among all unit vectors $v$ (unit for the chosen inner product), the
directional derivative $df_\theta(v)$ is maximal for $v = \nabla f(\theta)/\norm{\nabla f(\theta)}$, where it equals
$\norm{\nabla f(\theta)}$, and minimal for the opposite direction.
\end{theorem}

\begin{proof}
By Cauchy--Schwarz, $df_\theta(v) = \inner{\nabla f(\theta)}{v}\leq\norm{\nabla f(\theta)}\norm v = \norm{\nabla f(\theta)}$, with equality iff $v$ is a
positive multiple of $\nabla f(\theta)$.
\end{proof}

``Steepest'' is relative to the notion of unit step: changing the metric changes which direction
is steepest. This is the whole idea behind preconditioned methods
(\cref{gradient-tangent-space:sec:metric}).

\section{Level sets}\label{gradient-tangent-space:sec:level-sets}

\begin{proposition}[The gradient is orthogonal to level sets]\label{gradient-tangent-space:prop:level}
Let $\gamma$ be a differentiable curve with $\gamma(0) = \theta$ that stays in the level set
$\{f = f(\theta)\}$. Then $\inner{\nabla f(\theta)}{\gamma'(0)} = 0$.
\end{proposition}

\begin{proof}
$f(\gamma(t))$ is constant, so its derivative at $0$ vanishes: $0 = df_\theta(\gamma'(0)) = \inner{\nabla f(\theta)}{\gamma'(0)}$.
\end{proof}

So in parameter space the gradient crosses the contour lines at right angles, and gradient
descent steps across them (\cref{gradient-tangent-space:fig:contours}).

\begin{figure}[ht]
\centering
\begin{tikzpicture}[>=Stealth]
\begin{axis}[width=8.5cm, height=6cm, axis equal image, xmin=-3.3, xmax=3.3, ymin=-1.8, ymax=1.8,
  xlabel={$\theta_1$}, ylabel={$\theta_2$}, axis lines=middle, xtick=\empty, ytick=\empty]
  \foreach \c in {0.5,1.5,3,5,7.25,10}{
    \addplot[thin, bookaccent!70, domain=0:360, samples=80] ({sqrt(\c)*cos(x)}, {sqrt(\c)/2*sin(x)});
  }
  \addplot[thick, mark=*, mark size=1.4pt, bookpitfall] coordinates
    {(2.5,1.0) (2.0,0.2) (1.6,0.04) (1.28,0.008) (1.024,0.0016) (0.8192,0.0003) (0.6554,0.0001)};
  \draw[->, very thick, bookexample] (axis cs:2.5,1.0) -- (axis cs:2.5+0.1*5,1.0+0.1*8);
  \node[bookexample, right] at (axis cs:3.0,1.75) {$\nabla L$};
  \node[above left] at (axis cs:2.5,1.0) {$\theta_0$};
\end{axis}
\end{tikzpicture}
\caption{Level curves of $L(\theta) = \theta_1^2 + 4\theta_2^2$ in parameter space. The gradient at
$\theta_0 = (2.5,1)$ (green) is orthogonal to the level curve through $\theta_0$. Gradient descent with
$\eta = 0.1$ (red) multiplies $\theta_1$ by $0.8$ and $\theta_2$ by $0.2$ at each step: it first kills the steep
direction, then crawls along the flat one.}
\label{gradient-tangent-space:fig:contours}
\end{figure}

\section{The gradient lives in parameter space}\label{gradient-tangent-space:sec:graph}

The loss $L:\R^n\to\R$ has a graph $\Gamma = \{(\theta, L(\theta))\}\subset\R^{n+1}$. For $n = 2$ we can draw it as a
surface over the parameter plane, and it is tempting to draw the gradient as an arrow on that
surface, pointing downhill. That picture is wrong.

\begin{pitfall}[The gradient does not slide down the surface]
$\nabla L(\theta)\in\R^n$ has as many components as $\theta$; the graph lives in $\R^{n+1}$. The gradient is an
arrow in the \emph{parameter plane} (the floor of the 3D picture), orthogonal to the contour lines
there; gradient descent moves the point $\theta$ in that plane, and the loss value merely
follows. An arrow drawn on the surface is a tangent vector of the graph, with $n+1$ components ---
a different object (\cref{gradient-tangent-space:fig:surface}).
\end{pitfall}

\begin{figure}[ht]
\centering
\begin{tikzpicture}[>=Stealth]
\begin{axis}[width=9cm, height=7cm, view={-30}{30}, xmin=-2.6, xmax=2.6, ymin=-1.6, ymax=1.6, zmin=-4, zmax=8,
  xlabel={$\theta_1$}, ylabel={$\theta_2$}, zlabel={$L$}, xtick=\empty, ytick=\empty, ztick=\empty]
  % the floor: a copy of the parameter plane R^2, drawn below the graph
  \addplot3[fill=gray!12, draw=gray!50, thin] coordinates {(-2.6,-1.6,-4) (2.6,-1.6,-4) (2.6,1.6,-4) (-2.6,1.6,-4) (-2.6,-1.6,-4)};
  \node[gray, anchor=south west, font=\small] at (axis cs:-2.5,-1.5,-4) {$\R^2$};
  % level curves in the floor; the middle one passes through theta (L = 3.4)
  \addplot3[thin, bookaccent, domain=0:360, samples=40, samples y=1] ({1.0*cos(x)}, {0.5*sin(x)}, -4);
  \addplot3[thin, bookaccent, domain=0:360, samples=40, samples y=1] ({1.8439*cos(x)}, {0.92195*sin(x)}, -4);
  \addplot3[thin, bookaccent, domain=0:360, samples=40, samples y=1] ({2.4495*cos(x)}, {1.22474*sin(x)}, -4);
  % the graph of L
  \addplot3[surf, opacity=0.45, faceted color=bookaccent!40, colormap={flat}{color=(bookaccent!15) color=(bookaccent!15)}, domain=-2.2:2.2, y domain=-1.1:1.1, samples=17, samples y=13] {x^2 + 4*y^2};
  \addplot3[only marks, mark=*, mark size=1.6pt] coordinates {(1.8,0.2,3.4)};
  \addplot3[dashed, gray] coordinates {(1.8,0.2,-4) (1.8,0.2,3.4)};
  % gradient (3.6, 1.6) at theta = (1.8, 0.2), scaled by 0.2, drawn in the floor
  \draw[->, very thick, bookexample] (axis cs:1.8,0.2,-4) -- (axis cs:2.52,0.52,-4);
  \addplot3[only marks, mark=*, mark size=1.6pt, bookpitfall] coordinates {(1.8,0.2,-4)};
  \node[bookexample, anchor=south west] at (axis cs:2.52,0.52,-4) {$\nabla L(\theta)$};
  \node[bookpitfall, anchor=north] at (axis cs:1.8,0.2,-4) {$\theta$};
  \node[anchor=west] at (axis cs:1.8,0.2,3.4) {$(\theta, L(\theta))$};
\end{axis}
\end{tikzpicture}
\caption{The graph of $L$ lives in $\R^3$, but the parameter $\theta$ and the gradient $\nabla L(\theta)$ live in the
floor $\R^2$ (drawn below the graph for readability), where the level curves are. The gradient is
orthogonal to the level curve through $\theta$; it is not tangent to the surface.}
\label{gradient-tangent-space:fig:surface}
\end{figure}

If one insists on relating the gradient to the graph, the precise statement is the following.

\begin{proposition}[Gradient and the tangent space of the graph]\label{gradient-tangent-space:prop:graph}
The tangent vectors of $\Gamma$ at $(\theta_0, L(\theta_0))$ are the vectors $(v, \nabla L(\theta_0)^\top v)$, $v\in\R^n$. The
orthogonal projection of the vertical unit vector $e_{n+1}$ onto this tangent space --- the steepest
ascent direction \emph{on the surface} --- is
\[
  \frac1{1 + \norm{\nabla L(\theta_0)}^2}\begin{bmatrix}\nabla L(\theta_0)\\ \norm{\nabla L(\theta_0)}^2\end{bmatrix},
\]
whose first $n$ components are proportional to $\nabla L(\theta_0)$.
\end{proposition}

\begin{proof}
A curve $t\mapsto(\theta_0 + tv, L(\theta_0 + tv))$ on $\Gamma$ has velocity $(v, \nabla L^\top v)$ at $t = 0$. The tangent
space is the hyperplane orthogonal to $\mathbf n = (\nabla L, -1)$. Projecting $e_{n+1}$:
$e_{n+1} - \frac{\inner{e_{n+1}}{\mathbf n}}{\norm{\mathbf n}^2}\mathbf n = e_{n+1} + \frac{1}{1 + \norm{\nabla L}^2}(\nabla L, -1)$, which is the stated
vector.
\end{proof}

So the downhill arrow on the surface \emph{projects} onto a multiple of $-\nabla L$ in the floor, which is
what popular pictures are trying to show. For $n$ in the millions there is no surface to look at,
and only the parameter-space view remains.

\subsection{The hidden identification}\label{gradient-tangent-space:sec:affine}

The update $\theta\leftarrow\theta - \eta\nabla L(\theta)$ involves objects of two different kinds: $\theta$ is a
\emph{point} (a state of the model), while $\nabla L(\theta)$ and $-\eta\nabla L(\theta)$ are \emph{vectors} (a velocity at
$\theta$, and a displacement). The formula adds a displacement to a point. When parameters are
stored as arrays of $n$ numbers, points and vectors are both elements of $\R^n$ and the
distinction is invisible; it is worth making explicit, because the structure that legitimates
the update is not the one usually invoked.

\paragraph{Affine spaces.}

\begin{definition}[Affine space]\label{gradient-tangent-space:def:affine}
An \emph{affine space} is a set $\mathcal A$ (the \emph{points}) together with a real vector space $V$ (the
\emph{displacements}, or translations) and an operation $\mathcal A\times V\to\mathcal A$, written $(a, v)\mapsto a + v$,
such that for all $a, b\in\mathcal A$ and $u, v\in V$:
\begin{enumerate}[nosep]
  \item $a + 0 = a$;
  \item $(a + u) + v = a + (u + v)$;
  \item there is exactly one $v\in V$ with $a + v = b$; it is written $b - a$.
\end{enumerate}
The dimension of $\mathcal A$ is that of $V$.
\end{definition}

Point $+$ vector is a point and point $-$ point is a vector; that is all. There is no sum of two
points, no multiple of a point, and no distinguished point. Choosing a point $o\in\mathcal A$ as
\emph{origin} gives a bijection $\mathcal A\to V$, $a\mapsto a - o$, and choosing in addition a basis of $V$ gives
coordinates in $\R^n$. Both choices are arbitrary, like the zero of a temperature scale: the zero
of the Celsius scale is a point, whereas a difference of zero degrees is a vector, and it means the
same in every scale (an example of Hubbard and Hubbard). A vector space $V$ is an affine space over itself, with $a + v$ the vector
addition; this is how $\R^n$ is usually used, origin included.

Textbooks differ in how openly they say so. Hubbard and Hubbard (\emph{Vector Calculus, Linear
Algebra, and Differential Forms}, 5th ed., §1.1) separate ``positional data'' (points) from
``incremental data'' (vectors), give them different notations, and state that points cannot be
added while vectors can --- the axioms above without the name. Loomis and Sternberg
(\emph{Advanced Calculus}, §1.2) build the coordinate correspondence between Euclidean space and
$\R^3$ from a chosen origin, then deliberately transfer the geometric language to $\R^3$, where a
translation becomes the addition of a constant vector; from then on, as in
their definition of the differential (§3.6), a point and an increment are both elements of the
same normed vector space. Both are legitimate. The second is what our arrays do.

\begin{proposition}[Affine combinations]\label{gradient-tangent-space:prop:affine-comb}
Let $a_1,\dots,a_k\in\mathcal A$, $\lambda_1,\dots,\lambda_k\in\R$ and $o\in\mathcal A$.
\begin{enumerate}[nosep]
  \item If $\sum_i\lambda_i = 1$, the point $o + \sum_i\lambda_i(a_i - o)$ does not depend on $o$. It is written
        $\sum_i\lambda_ia_i$ (an \emph{affine combination}, or barycentre).
  \item If $\sum_i\lambda_i = 0$, the vector $\sum_i\lambda_i(a_i - o)$ does not depend on $o$.
  \item Otherwise both depend on $o$.
\end{enumerate}
\end{proposition}

\begin{proof}
Replace $o$ by $o' = o + w$. Then $a_i - o' = (a_i - o) - w$, so
$\sum_i\lambda_i(a_i - o') = \sum_i\lambda_i(a_i - o) - \bigl(\sum_i\lambda_i\bigr)w$ and
\[
  o' + \sum_i\lambda_i(a_i - o') = o + \sum_i\lambda_i(a_i - o) + \Bigl(1 - \sum_i\lambda_i\Bigr)w .
\]
The vector is unchanged for every $w$ iff $\sum_i\lambda_i = 0$, and the point iff $\sum_i\lambda_i = 1$.
\end{proof}

The midpoint $\frac12a + \frac12b$ and the mean of a cloud of points are affine combinations; $a + b$ and
$2a$ are not, and $b - a$ (coefficients $1$ and $-1$) is a vector.

\paragraph{Tangent vectors and the identification.} Let $\gamma$ be a curve in $\mathcal A$ with $\gamma(0) = \theta$.
For each $t$, $\gamma(t) - \theta$ is a vector of $V$, so the velocity
\[
  \gamma'(0) = \lim_{t\to0}\frac{\gamma(t) - \theta}{t}
\]
is a limit in the finite-dimensional space $V$ (whose topology does not depend on the norm). The
tangent space at $\theta$ --- the set of velocities of curves through $\theta$ --- is therefore $V$ itself,
for every $\theta$: the curve $t\mapsto\theta + tv$ has velocity $v$, and every velocity arises this way. This
identification of each tangent space $T_\theta\mathcal A$ with the one space $V$ involves no choice; it is the
identification that $\theta - \eta\nabla L(\theta)$ uses without saying so. Already for $f:\R\to\R$, the derivative
$f'(a)$ is a velocity at $a$ that we read as a displacement along the axis.

\paragraph{Three structures, three objects.} Each object in the update needs one more piece of
structure than the previous one.
\begin{itemize}[nosep]
  \item \emph{Affine structure} $(\mathcal A, V)$. Points, displacements, and the identification
        $T_\theta\mathcal A = V$. Enough to write $\theta + v$ and to define velocities.
  \item \emph{The dual space} $V^*$, which needs nothing more. For $L:\mathcal A\to\R$, the differential is the linear
        form $dL_\theta\in V^*$ with $L(\theta + v) = L(\theta) + dL_\theta(v) + o(\norm v)$
        (\cref{gradient-tangent-space:sec:differential}); like the velocity, it involves no inner
        product and no origin.
  \item \emph{An inner product on $V$.} It turns the form $dL_\theta$ into the vector $\nabla L(\theta)\in V$, by
        $\inner{\nabla L(\theta)}{v} = dL_\theta(v)$. Different inner products give different gradients
        (\cref{gradient-tangent-space:prop:metric}).
\end{itemize}
An affine space whose space of displacements carries an inner product is a \emph{Euclidean affine
space}. This, not $\R^n$ with its origin and its standard basis, is the structure that gradient
descent requires: the origin is never used, and the basis enters only through the inner
product.

\begin{pitfall}[``$\R^n$ is flat'' is true, but it is not the reason]
Flatness is a property of a Riemannian metric (zero curvature). It is neither necessary nor
sufficient for $\theta + v$ to make sense.
\begin{itemize}[nosep]
  \item An affine space without inner product has no curvature at all, flat or not, and $\theta + v$
        is defined.
  \item A circle of length $1$, $\R/\Z$, is flat (every one-dimensional Riemannian manifold is), yet
        it is not an affine space: the displacement from $0.1$ to $0.3$ could be $0.2$, $1.2$ or $-0.8$
        (\cref{gradient-tangent-space:exr:6}). Flatness gives the affine picture only locally.
\end{itemize}
What the update uses is the affine structure; the inner product then supplies the gradient.
\end{pitfall}

\paragraph{A test for hidden origins.} Since the origin is arbitrary, a formula has a geometric
meaning only if it gives the same result whichever origin is used to compute it. Concretely, move
the origin by $w$: a point with coordinates $x$ gets coordinates $x' = x - w$, a vector keeps its
coordinates, and the loss becomes $L'(x') = L(x' + w)$.

\begin{proposition}[Gradient descent ignores the origin]\label{gradient-tangent-space:prop:origin}
$\nabla L'(x') = \nabla L(x' + w)$. Consequently, if $x_{k+1} = x_k - \eta\nabla L(x_k)$ and $x'_k = x_k - w$, then
$x'_{k+1} = x'_k - \eta\nabla L'(x'_k)$: the two computations describe the same points.
\end{proposition}

\begin{proof}
A translation has the identity as differential, so the chain rule gives $\nabla L'(x') = \nabla L(x' + w)$.
Then $x'_{k+1} = x_{k+1} - w = (x_k - w) - \eta\nabla L(x_k) = x'_k - \eta\nabla L'(x'_k)$.
\end{proof}

The test is useful because the computer does not apply it: on arrays, $x + y$ and $2x$ are as
easy to write as $x + v$. Here is how common updates fare.
\begin{itemize}[nosep]
  \item \emph{Momentum} (heavy ball), $\theta_{k+1} = \theta_k - \eta\nabla L(\theta_k) + \beta(\theta_k - \theta_{k-1})$: the added term is a
        difference of points, a vector. It passes.
  \item \emph{Iterate averaging}, $\bar\theta_K = \frac1K\sum_{k=1}^K\theta_k$ (Polyak averaging), and the exponential moving
        average $\bar\theta\leftarrow\beta\bar\theta + (1 - \beta)\theta$: affine combinations, with weights summing to $1$. They pass.
  \item \emph{Weight decay}, $\theta\leftarrow\theta - \eta\bigl(\nabla L(\theta) + \lambda\theta\bigr) = (1 - \eta\lambda)\theta - \eta\nabla L(\theta)$: the
        coefficients of $\theta$ do not sum to $1$. It fails (\cref{gradient-tangent-space:ex:decay}).
\end{itemize}

\begin{example}[Weight decay has a hidden origin]\label{gradient-tangent-space:ex:decay}
Written honestly, the term $\lambda\theta$ is $\lambda(\theta - o)$, where $o$ is the point whose coordinates are $0$:
\[
  \theta\leftarrow\theta - \eta\bigl(\nabla L(\theta) + \lambda(\theta - o)\bigr)
  = \bigl(o + (1 - \eta\lambda)(\theta - o)\bigr) - \eta\nabla L(\theta),
\]
a contraction of ratio $1 - \eta\lambda$ towards $o$, followed by a gradient step. It is gradient descent
on $L(\theta) + \frac\lambda2\norm{\theta - o}^2$, that is, MAP estimation with a Gaussian prior centred at $o$
(\cref{bayesian-reasoning:sec:map}). The test does not say that weight decay is wrong; it says that
the formula contains a modelling choice --- the centre of the prior --- that the notation $\lambda\theta$ hides.
For a network trained from scratch, decaying towards the zero weights is a reasonable prior. For a
network fine-tuned from pre-trained weights $\theta_0$, decaying towards $\theta_0$, i.e.\ replacing $\lambda\theta$ by
$\lambda(\theta - \theta_0)$, is often the better one (the ``L2-SP'' regularizer of Li, Grandvalet and Davoine,
2018).
\end{example}

The same test applied to a change of \emph{basis} rather than of origin is the subject of
\cref{gradient-tangent-space:exr:2}. Plain gradient descent is unchanged by an orthonormal change
of basis but not by a general one, because ``the dot product of the coordinates'' is a different
inner product in each non-orthonormal basis. Coordinate-wise methods such as Adam pass the origin
test but are changed even by a rotation of an orthonormal basis
(\cref{gradient-tangent-space:exr:4}): they use the basis itself, not only an inner product.

On a curved parameter space there is no space of displacements at all: the tangent spaces at
different points are different vector spaces, $\theta + v$ is not defined, and the update must be
rethought (\cref{gradient-tangent-space:sec:manifolds}).

\section{Changing the metric}\label{gradient-tangent-space:sec:metric}

By \cref{gradient-tangent-space:prop:metric}, gradient descent for the metric $M$ is
\[
  \theta\leftarrow\theta - \eta\,M^{-1}\nabla L(\theta):
\]
\emph{preconditioned} gradient descent is ordinary gradient descent for another inner product.
Three choices of $M$ are classical.
\begin{itemize}
  \item \textbf{Diagonal rescaling}, $M = \diag(m_1,\dots,m_n)$: one learning rate per coordinate. Adaptive
        methods (AdaGrad, RMSProp, Adam) estimate such a diagonal from past gradients
        (Part~VI).
  \item \textbf{Newton's method}, $M = \nabla^2L(\theta)$ (when positive definite): the gradient for the metric
        defined by the local curvature.
  \item \textbf{Natural gradient}, $M = F(\theta)$, the Fisher information matrix of a probabilistic model:
        the metric that measures parameter changes by the change they cause in the model's
        distribution (locally, the KL divergence), independently of the parametrization.
\end{itemize}

\begin{example}[The right metric removes the zig-zag]
For $L(\theta) = \theta_1^2 + 4\theta_2^2$, the Euclidean gradient is $(2\theta_1, 8\theta_2)$ and gradient descent crawls
(\cref{gradient-tangent-space:fig:contours}). With $M = \nabla^2L = \diag(2, 8)$,
$M^{-1}\nabla L(\theta) = (\theta_1, \theta_2) = \theta$: the gradient for this metric points straight at the minimum, and
one step with $\eta = 1$ reaches it. In the metric $M$ the level sets are circles.
\end{example}

\section{Gradients on manifolds}\label{gradient-tangent-space:sec:manifolds}

Some parameters are constrained to curved sets: unit vectors (the sphere), orthonormal frames
(the Stiefel manifold $\{\Theta\in\R^{n\times p} : \Theta^\top\Theta = I\}$), rotations, probability vectors.

\begin{definition}[Tangent space]
Let $\mathcal M\subset\R^d$ be a smooth surface (submanifold) and $\theta\in\mathcal M$. The tangent space $T_\theta\mathcal M$ is the set
of velocities $\gamma'(0)$ of smooth curves $\gamma$ that stay in $\mathcal M$ with $\gamma(0) = \theta$.
\end{definition}

$T_\theta\mathcal M$ is a vector space: the directions in which one can move \emph{while staying on} $\mathcal M$. A
\emph{Riemannian metric} assigns an inner product $g_\theta$ to each tangent space; on a submanifold of
$\R^d$ the natural choice is the restriction of the dot product. The Riemannian gradient
$\operatorname{grad}f(\theta)\in T_\theta\mathcal M$ is then defined, as before, by $g_\theta(\operatorname{grad}f(\theta), v) = df_\theta(v)$ for all
$v\in T_\theta\mathcal M$. It lies in the tangent space by construction, and the steepest-ascent theorem holds
within $T_\theta\mathcal M$.

\begin{proposition}[Riemannian gradient = projected Euclidean gradient]\label{gradient-tangent-space:prop:projection}
If $f$ is the restriction to $\mathcal M$ of a function $F$ defined on $\R^d$, and $\mathcal M$ carries the induced
metric, then $\operatorname{grad}f(\theta) = P_\theta\nabla F(\theta)$, the orthogonal projection of the Euclidean gradient onto
$T_\theta\mathcal M$.
\end{proposition}

\begin{proof}
For $v\in T_\theta\mathcal M$, $df_\theta(v) = dF_\theta(v) = \inner{\nabla F(\theta)}{v} = \inner{P_\theta\nabla F(\theta)}{v}$, because
$\nabla F - P_\theta\nabla F$ is orthogonal to $T_\theta\mathcal M$. Since $P_\theta\nabla F(\theta)\in T_\theta\mathcal M$, it satisfies the defining
property of the Riemannian gradient.
\end{proof}

So a projection is needed only when the function was defined on the surrounding space: the
\emph{ambient} gradient $\nabla F$ may point off the surface, and projecting it gives the intrinsic one.
An update $\theta - \eta\operatorname{grad}f(\theta)$ still leaves $\mathcal M$ (a tangent step leaves a curved surface), so it is
followed by a \emph{retraction} $R_\theta$ back to $\mathcal M$:
$\theta\leftarrow R_\theta\bigl(-\eta\operatorname{grad}f(\theta)\bigr)$. In $\R^n$, $R_\theta(v) = \theta + v$ and we recover gradient descent.

\begin{example}[The sphere and PCA]\label{gradient-tangent-space:ex:sphere}
On the unit sphere $S^{d-1} = \{\theta : \norm\theta = 1\}$, curves with $\norm{\gamma(t)}^2 = 1$ satisfy
$2\theta^\top\gamma'(0) = 0$, so $T_\theta S^{d-1} = \{v : \theta^\top v = 0\}$ and $P_\theta = I - \theta\theta^\top$. A simple retraction is
normalization, $R_\theta(v) = (\theta + v)/\norm{\theta + v}$.

Take $F(\theta) = \theta^\top S\theta$ with $S$ symmetric (a covariance matrix). Then $\nabla F = 2S\theta$ and
\[
  \operatorname{grad}f(\theta) = (I - \theta\theta^\top)\,2S\theta = 2\bigl(S\theta - (\theta^\top S\theta)\,\theta\bigr).
\]
It vanishes iff $S\theta$ is a multiple of $\theta$: the critical points on the sphere are the unit eigenvectors
of $S$, and the maximum of the variance $\theta^\top S\theta$ is reached at the leading one --- the first principal
axis of \cref{ch:pca-overview}. Riemannian gradient ascent with the normalization retraction is
a smoothed version of the power method.
\end{example}

\begin{figure}[ht]
\centering
\begin{tikzpicture}[>=Stealth, scale=1.5]
  \draw[thick, bookaccent] (0,0) circle (1);
  \coordinate (T) at (0.6,0.8);
  \draw[gray] ($(T)+(-1.2,0.9)$) -- ($(T)+(1.2,-0.9)$) node[right, black] {$T_\theta S^1$};
  \fill (T) circle (1.3pt) node[above left] {$\theta$};
  \draw[->, thick] (T) -- ($(T)+(1.0,0.5)$) node[right] {$\nabla F(\theta)$};
  \draw[->, very thick, bookexample] (T) -- ($(T)+(0.4,-0.3)$) node[below right] {$\operatorname{grad}f(\theta)$};
  \draw[dashed] ($(T)+(1.0,0.5)$) -- ($(T)+(0.4,-0.3)$);
  \fill (0,0) circle (1pt);
\end{tikzpicture}
\caption{On a circle, the ambient gradient $\nabla F(\theta)$ is split into a normal part (which would push
$\theta$ off the circle) and a tangent part, the Riemannian gradient.}
\label{gradient-tangent-space:fig:sphere}
\end{figure}

\begin{chaptersummary}
\begin{itemize}[nosep]
  \item The differential $df_\theta$ is a linear form; the gradient is the vector representing it for a chosen
        inner product: $\nabla f$ for the dot product, $M^{-1}\nabla f$ for $u^\top Mv$.
  \item The gradient is the steepest-ascent direction for the chosen notion of unit step, and it is
        orthogonal to the level sets.
  \item The gradient lives in parameter space, not on the graph of the loss; the downhill arrow on
        the surface only projects onto a multiple of $-\nabla L$.
  \item The update adds a displacement to a point: it needs an affine structure (for $\theta + v$ and
        velocities), nothing more for the differential, and an inner product for the gradient. It
        never uses the origin, whereas weight decay does: it is a prior centred at the zero weights.
  \item Preconditioning, Newton's method and the natural gradient are gradient descent in other
        metrics.
  \item On a manifold, the Riemannian gradient lies in the tangent space; for an embedded surface it
        is the projection of the ambient gradient, and steps are followed by a retraction. On the sphere,
        maximizing $\theta^\top S\theta$ recovers the first principal axis.
\end{itemize}
\end{chaptersummary}

\section*{Exercises}
\addcontentsline{toc}{section}{Exercises}

\begin{exercise}\label{gradient-tangent-space:exr:1}
Let $f(\theta) = \theta_1\theta_2$. Compute its gradient for the dot product and for $M = \diag(1, 4)$, at
$\theta = (1,1)$, and the corresponding unit steepest-ascent directions.
\solutionlink{gradient-tangent-space:exr:1}
\end{exercise}

\begin{exercise}\label{gradient-tangent-space:exr:2}
Show that if $\theta = A\phi$ with $A$ invertible (a reparametrization), Euclidean gradient descent on $\phi$
corresponds to gradient descent on $\theta$ for the metric $M = (AA^\top)^{-1}$. Conclude that plain gradient
descent depends on the parametrization.
\solutionlink{gradient-tangent-space:exr:2}
\end{exercise}

\begin{exercise}\label{gradient-tangent-space:exr:3}
For $S = \diag(3, 1)$ and $\theta = (\cos\alpha, \sin\alpha)$ on the unit circle, compute $\operatorname{grad}f(\theta)$ for
$f(\theta) = \theta^\top S\theta$, and check that it vanishes exactly at the four points $(\pm1,0)$, $(0,\pm1)$.
\solutionlink{gradient-tangent-space:exr:3}
\end{exercise}

\begin{exercise}\label{gradient-tangent-space:exr:4}
Adam updates $m\leftarrow\beta_1m + (1 - \beta_1)g$, $s\leftarrow\beta_2s + (1 - \beta_2)\,g\odot g$ from the gradient $g = \nabla L(\theta)$ and steps
$\theta\leftarrow\theta - \eta\,\hat m\oslash(\sqrt{\hat s} + \varepsilon)$, coordinate-wise, with the bias-corrected $\hat m$, $\hat s$.
\begin{enumerate}[nosep]
  \item Show that Adam passes the origin test of \cref{gradient-tangent-space:prop:origin}.
  \item Show that its first step (from $m = s = 0$) is $-\eta\,\mathrm{sign}(g)$ coordinate-wise when $\varepsilon = 0$
        (with $\mathrm{sign}(0) = 0$).
  \item In $\R^2$ let $g = (2, 1)$. Compute the first step in the standard basis, then in the orthonormal
        basis rotated by $45^\circ$, and express the second back in the standard basis. Are the two steps
        the same? Parallel?
\end{enumerate}
\solutionlink{gradient-tangent-space:exr:4}
\end{exercise}

\begin{exercise}\label{gradient-tangent-space:exr:5}
Run weight decay, $x\leftarrow(1 - \eta\lambda)x - \eta\nabla L(x)$, in the coordinates $x$ and in the coordinates $x' = x - w$
of a shifted origin, i.e.\ $x'\leftarrow(1 - \eta\lambda)x' - \eta\nabla L'(x')$ with $L'(x') = L(x' + w)$. Starting from the same
point, by how much do the two new points differ after one step? For $L\equiv0$, where do the two
sequences converge?
\solutionlink{gradient-tangent-space:exr:5}
\end{exercise}

\begin{exercise}\label{gradient-tangent-space:exr:6}
On the circle $\mathcal C = \R/\Z$, whose points are the classes $[x] = x + \Z$, let $v\in\R$ act by $[x] + v = [x + v]$.
Show that this is well defined and satisfies axioms 1 and 2 of
\cref{gradient-tangent-space:def:affine}, but not axiom 3. Is $\theta - \eta\nabla f(\theta)$ still meaningful for a
differentiable $f$ on $\mathcal C$?
\solutionlink{gradient-tangent-space:exr:6}
\end{exercise}
\end{notechapter}
